\documentclass[11pt]{article}
\usepackage[a4paper,margin=28mm]{geometry}
\usepackage{amsmath,amsthm}
\usepackage[hidelinks]{hyperref}
\newtheorem{theorem}{Theorem}
\title{Subset Sum by item count under the Satisfiability Hypothesis}
\author{Atlas\thanks{The proof in this note was found using ChatGPT.}}
\date{18 November 2026}
\begin{document}
\maketitle

An instance of \textsc{Subset Sum} is an indexed list of \(p\) positive
integers and a positive target. All numbers are in binary. The question is
whether some subset of the indices has values summing to the target.
Let \(B\) be the complete input length in bits.

\hypertarget{subset-sum-item-count-sh}{}
\begin{theorem}\label{thm:subset-sum-sh}
Assume the deterministic decision Satisfiability Hypothesis (SH) for
CNF satisfiability by normalized input length. For each fixed
\(0<\varepsilon<1\), there is no deterministic
\(2^{O(p^{1-\varepsilon})}B^{O(1)}\)-time algorithm for \textsc{Subset Sum}.
\end{theorem}

\begin{proof}
We use the consequence of SH for width-three CNF by the combined number
of variables and clauses~\cite[Lemma~2.1]{sh}. The length measured in
bits after renumbering occurring variables introduces at most a logarithmic
factor in that measure, which preserves the existence of a fixed power
saving in the exponent.

Let the formula have \(v\) variables and \(m\) clauses, each listing at most
three literal occurrences. Formulas with no clauses or with an empty clause
can be decided directly. Use one base-eight digit position per variable and
per clause. For each variable create two items, one for each assignment to
that variable. Each has digit one in that variable's position and zero in
the other variable positions. In a clause position its digit is the number
of literal occurrences made true by the corresponding assignment to that
variable. Create two additional indexed items per clause. Each has digit
one in its own clause position and zeros elsewhere. Set the target digits
to one at the variable positions and three at the clause positions.
These are positive integers. The construction is a variant of the
standard digit reduction~\cite[Section~3]{weighted}.

There is no carry in any subset sum. A variable position has total at
most two and a clause position has total at most five, since even when
all assignment items are selected, each of the at most three occurrences
contributes once. Equality to the target forces the choice of exactly one
assignment item per variable. In a clause, the two extra items allow that
assignment's digit sum to reach three precisely when at least one
occurrence is true. Thus the formula is satisfiable exactly when the
constructed instance is a yes-instance.

There are \(p=2v+2m\) items and all numbers have at most \(3(v+m)\)
bits, so construction and output length are polynomial in the complete
formula input length. For any fixed \(0<\varepsilon<1\), the stated
algorithm would solve width-three satisfiability in time
\(2^{O((v+m)^{1-\varepsilon})}\) times a polynomial in that input length.
This contradicts the cited consequence of SH.
\end{proof}

\begin{thebibliography}{9}
\bibitem{sh}
MIT Hardness Group, J.~Brunner, E.~D. Demaine, F.~Frei, T.~Gomez,
M.~Hecher, J.~Li, and J.~Lynch.
\emph{ETH is Dead, Long Live SH}.
\href{https://arxiv.org/abs/2611.02650v1}{arXiv:2611.02650v1}, 2026.

\bibitem{weighted}
M.~Etscheid, S.~Kratsch, M.~Mnich, and H.~Röglin.
Polynomial kernels for weighted problems.
\emph{Journal of Computer and System Sciences}, 84:1--10, 2017.
\href{https://doi.org/10.1016/j.jcss.2016.06.004}{DOI}.
\end{thebibliography}
\end{document}
